Large language models are remarkably data-hungry, yet humans generalize from far fewer examples. We hypothesize that rapid memorization is a primary bottleneck: neural networks collapse training examples into lookup tables before extracting generalizable structure, and this premature memorization limits data efficiency. We test this hypothesis using modular arithmetic, a controlled setting where memorization and generalization are cleanly separable via the grokking phenomenon. Across three experiments, we systematically vary forgetting pressure (weight decay), data redundancy (duplication), and forgetting mechanisms (dropout, noise injection, weight reset) while tracking per-example forgetting events. Our key finding is that increasing forgetting pressure via weight decay accelerates generalization by 30$\times$ (from ${\sim}$17,850 steps to ${\sim}$583 steps), with a strong negative correlation between forgetting pressure and time-to-generalization (Spearman $\rho = -0.88$, $p < 0.00002$). We further show that data duplication slows generalization by 37\%, and that an optimal level of forgetting exists---too little leads to prolonged memorization while too much prevents learning entirely. These results provide direct evidence that models which forget individual examples more readily discover generalizable solutions faster, reframing data efficiency as a forgetting problem rather than a data quantity problem.
